% 原PDF165页附录4.G标题—171页附录4.H标题前。
\section*{附录4.G\quad 动态结构方程模型的\texorpdfstring{$L_1$}{L1}惩罚最大似然估计的近端算法}
\addcontentsline{toc}{section}{附录4.G 动态结构方程模型的L1惩罚最大似然估计的近端算法}
\markright{附录4.G L1惩罚最大似然估计的近端算法}

令
\begin{equation}
\begin{aligned}
f(\Gamma,A,B)&=-l(\Gamma,A,B,\Sigma\mid Y)\\
&=\frac{nTm}{2}\log(2\pi)+\frac n2\log|\Sigma|-n\log|\Gamma|
+\frac12\sum_{i=1}^{n}e_i'\Sigma^{-1}e_i,
\end{aligned}\tag{4.G.1}
\end{equation}
其中
\[
\Gamma=\begin{bmatrix}\Gamma^1&0&\cdots&0\\0&\Gamma^2&\cdots&0\\\vdots&\vdots&\vdots&\vdots\\0&0&\cdots&\Gamma^T\end{bmatrix},\quad
A=\begin{bmatrix}A^1&0&\cdots&0\\0&A^2&\cdots&0\\\vdots&\vdots&\vdots&\vdots\\0&0&\cdots&A^T\end{bmatrix},
\]
\[
B=\begin{bmatrix}B^1&0&\cdots&0\\0&B^2&\cdots&0\\\vdots&\vdots&\vdots&\vdots\\0&0&\cdots&B^T\end{bmatrix},\quad
y_i=\begin{bmatrix}y_i^1\\y_i^2\\\vdots\\y_i^T\end{bmatrix},\quad
y_i^t=\begin{bmatrix}y_{i1}^t\\y_{i2}^t\\\vdots\\y_{iM}^t\end{bmatrix},
\]
\[
y_{li}=\begin{bmatrix}y_{li}^0\\y_{li}^1\\\vdots\\y_{li}^{T-1}\end{bmatrix},\quad
y_{li}^{t-1}=\begin{bmatrix}y_{li1}^{t-1}\\y_{li2}^{t-1}\\\vdots\\y_{liM}^{t-1}\end{bmatrix},\quad
X_i=\begin{bmatrix}X_i\\X_i\\\vdots\\X_i\end{bmatrix},\quad
x_i=\begin{bmatrix}x_{i1}\\x_{i2}\\\vdots\\x_{iK}\end{bmatrix},
\]
\[
e_i=\begin{bmatrix}e_i^1\\e_i^2\\\vdots\\e_i^T\end{bmatrix},\quad
e_i^t=\begin{bmatrix}e_{i1}^t\\\vdots\\e_{iM}^t\end{bmatrix},\qquad
e_i=-(\Gamma'y_i+A'y_{li}+B'X_i).
\]
目标是最小化
\begin{equation}
\begin{aligned}
\min_{\Gamma^t,A^t,B^t}\quad
f(\Gamma,A,B)+\lambda_1\sum_{t=1}^{T}\|\Gamma^t\|_1
+\lambda_2\sum_{t=1}^{T}\|A^t\|_1
+\lambda_3\sum_{t=1}^{T}\|B^t\|_1.
\end{aligned}\tag{4.G.2}
\end{equation}
问题（4.G.2）的目标函数由凸可微函数和非光滑函数组成。传统牛顿法可以高效求解无约束光滑优化问题，但不适于大型非光滑凸问题。近端方法可视为牛顿法由光滑优化向非光滑优化的扩展。

将近端方法用于问题（4.G.2）的第一步，是对 $f(\Gamma,A,B)$ 关于 $\Gamma,A,B$ 求偏导。应用第1.4节的矩阵微积分，得到
\begin{equation}
\frac{\partial f}{\partial\Gamma'}=-n\Gamma^{-1}
-\sum_{i=1}^{n}(\Sigma^{-1})e_iy_i'.\tag{4.G.3}
\end{equation}
两侧取转置，得到
\begin{equation}
\frac{\partial f}{\partial\Gamma}=-n(\Gamma')^{-1}
-\sum_{i=1}^{n}y_ie_i'(\Sigma^{-1}).\tag{4.G.4}
\end{equation}
类似地，
\begin{equation}
\frac{\partial f}{\partial A}=-\sum_{i=1}^{n}y_{li}e_i'\Sigma^{-1}\tag{4.G.5}
\end{equation}
以及
\begin{equation}
\frac{\partial f}{\partial B}=-\sum_{i=1}^{n}X_ie_i'\Sigma^{-1},\tag{4.G.6}
\end{equation}
其中
\begin{equation}
e_i'=-(y_i'\Gamma+y_{li}'A+X_i'B).\tag{4.G.7}
\end{equation}
如果不同时间点的误差独立，则 $\Sigma=\operatorname{diag}(\Sigma_{11},\Sigma_{22},\ldots,\Sigma_{TT})$ 为分块对角矩阵。式（4.G.4）—（4.G.7）重写为
\begin{equation}
\frac{\partial f}{\partial\Gamma^t}=-n((\Gamma^t)')^{-1}
-\sum_{i=1}^{n}y_i^t(e_i^t)'\Sigma_{tt}^{-1},\tag{4.G.8}
\end{equation}
\begin{equation}
\frac{\partial f}{\partial A^t}=-\sum_{i=1}^{n}y_{li}^t(e_i^t)'\Sigma_{tt}^{-1},\tag{4.G.9}
\end{equation}
\begin{equation}
\frac{\partial f}{\partial B^t}=-\sum_{i=1}^{n}x_i(e_i^t)'\Sigma_{tt}^{-1},\tag{4.G.10}
\end{equation}
以及
\begin{equation}
(e_i^t)'=-((y_i^t)'\Gamma^t+(y_{li}^t)'A^t+x_i'B^t).\tag{4.G.11}
\end{equation}
令
\[
W_\Gamma=\operatorname{vec}(\Gamma),\quad
W_A=\operatorname{vec}(A),\quad W_B=\operatorname{vec}(B),
\]
\[
\frac{\partial f}{\partial W_\Gamma}=\operatorname{vec}\left(\frac{\partial f}{\partial\Gamma}\right),\quad
\frac{\partial f}{\partial W_A}=\operatorname{vec}\left(\frac{\partial f}{\partial A}\right),\quad
\frac{\partial f}{\partial W_B}=\operatorname{vec}\left(\frac{\partial f}{\partial B}\right),
\]
\begin{equation}
W=\begin{bmatrix}W_\Gamma\\W_A\\W_B\end{bmatrix},\qquad
\frac{\partial f}{\partial W}=\begin{bmatrix}
\dfrac{\partial f}{\partial W_\Gamma}\\
\dfrac{\partial f}{\partial W_A}\\
\dfrac{\partial f}{\partial W_B}
\end{bmatrix}.\tag{4.G.12}
\end{equation}
令
\[
f(W)=f(\Gamma,A,B),\quad
\Omega_\Gamma(W_\Gamma)=\|\Gamma\|_1,\quad
\Omega_A=\|A\|_1,\quad\Omega_B=\|B\|,
\]
以及
\[
\Omega(W)=\lambda_1\Omega_\Gamma(W_\Gamma)
+\lambda_2\Omega_A(W_A)+\lambda_3\Omega_B(W_B).
\]
优化问题（4.G.2）可写为
\begin{equation}
\min_W\quad f(W)+\Omega(W).\tag{4.G.13}
\end{equation}
回顾函数 $\Omega(w)$ 的近端算子（式（1.42））定义：
\begin{equation}
\operatorname{Prox}_\Omega(u)=\arg\min_{w\in\mathbb R^p}
\left(\Omega(W)+\frac12\|W-u\|_2^2\right).\tag{4.G.14}
\end{equation}
由近端算子的性质，
\begin{equation}
\operatorname{Prox}_\Omega(u)=\begin{bmatrix}
\operatorname{Prox}_{\lambda_1\Omega_\Gamma}(u_\Gamma)\\
\operatorname{Prox}_{\lambda_2\Omega_A}(u_A)\\
\operatorname{Prox}_{\lambda_3\Omega_B}(u_B)
\end{bmatrix},\tag{4.G.15}
\end{equation}
其中
\begin{equation}
u_\Gamma=W_\Gamma-\frac1L\frac{\partial f}{\partial W_\Gamma},\quad
u_A=W_A-\frac1L\frac{\partial f}{\partial W_A},\quad
u_B=W_B-\frac1L\frac{\partial f}{\partial W_B}.\tag{4.G.16}
\end{equation}
令 $u_{\Gamma ij},u_{Aij},u_{Bij}$ 分别为矩阵 $u_\Gamma,u_A,u_B$ 的第 $(i,j)$ 个元素。定义
\begin{equation}
Z=\begin{bmatrix}Z_\Gamma\\Z_A\\Z_B\end{bmatrix}
=\begin{bmatrix}
\operatorname{Prox}_{\lambda_1\Omega_\Gamma}(u_\Gamma)\\
\operatorname{Prox}_{\lambda_2\Omega_A}(u_A)\\
\operatorname{Prox}_{\lambda_3\Omega_B}(u_B)
\end{bmatrix}.\tag{4.G.17}
\end{equation}
令
\[
\begin{aligned}
u_{\Gamma ij}^k&=\gamma_{ij}^k-\lambda_1^k\left(\frac{\partial f}{\partial\Gamma}\right)_{ij},\\
u_{Aij}^k&=\alpha_{ij}^k-\lambda_2^k\left(\frac{\partial f}{\partial A}\right)_{ij},\qquad
u_{Bij}^k=\beta_{ij}^k-\lambda_3^k\left(\frac{\partial f}{\partial B}\right)_{ij}.
\end{aligned}
\]
非光滑函数 $\Omega_\Gamma,\Omega_A,\Omega_B$ 为 $L_1$ 范数。利用式（1.137）（Xiong 2018），得到三个近端算子的元素：
\begin{equation}
\gamma_{ij}^z=\operatorname{sign}(u_{\Gamma ij}^k)(|u_{\Gamma ij}^k|-\lambda_1^k)_+,
\quad i,j=1,\ldots,MT,\tag{4.G.18}
\end{equation}
\begin{equation}
\alpha_{ij}^z=\operatorname{sign}(u_{Aij}^k)(|u_{Aij}^k|-\lambda_2^k)_+,
\quad i,j=1,\ldots,MT,\tag{4.G.19}
\end{equation}
\begin{equation}
\beta_{ij}^z=\operatorname{sign}(u_{Bij}^k)(|u_{Bij}^k|-\lambda_3^k)_+,
\quad i=1,\ldots,KT,\ j=1,\ldots,MT,\tag{4.G.20}
\end{equation}
其中
\[
\Gamma_z=(\gamma_{ij}^z)_{MT\times MT},\quad
A_z=(\alpha_{ij}^z)_{MT\times MT},\quad B_z=(\beta_{ij})_{KT\times MT}.
\]
结合式（4.G.17）—（4.G.20），
\begin{equation}
Z_\Gamma=\operatorname{vec}(\Gamma_z),\quad
Z_A=\operatorname{vec}(A_z),\quad Z_B=\operatorname{vec}(B_z).\tag{4.G.21}
\end{equation}
令
\[
G_\Gamma^k=\operatorname{vec}(\Gamma^k),\quad
G_A^k=\operatorname{vec}(A^k),\quad G_B^k=\operatorname{vec}(B^k),\quad
G^k=\begin{bmatrix}G_\Gamma^k\\G_A^k\\G_B^k\end{bmatrix},
\]
其中
\[
\Gamma^k=(\gamma_{ij}^k)_{MT\times MT},\quad
A^k=(\alpha_{ij}^k)_{MT\times MT},\quad B^k=(\beta_{ij}^k)_{KT\times KT}.
\]
定义辅助函数
\begin{equation}
\begin{aligned}
f_\lambda(Z,G^k)={}&f(Z)
+\left(\frac{\partial f}{\partial Z}\right)^T(G^k-Z)
+\frac1{2\lambda_1}\|G_\Gamma^k-Z_\Gamma\|_2^2\\
&{}+\frac1{2\lambda_2}\|G_A^k-Z_A\|_2^2
+\frac1{2\lambda_3}\|G_B^k-Z_B\|_2^2.
\end{aligned}\tag{4.G.22}
\end{equation}

\begin{bookalgorithm}{4.G.1：近端梯度算法}
\textbf{步骤1：}给定 $W_\Gamma^k,W_A^k,W_B^k,G^k,G_\Gamma^k,G_A^k,G_B^k$、
$\lambda_1^{k-1},\lambda_2^{k-1},\lambda_3^{k-1}$，以及参数 $\alpha_1,\alpha_2,\alpha_3\in(0,1)$。
令 $\lambda_1=\lambda_1^{k-1}$、$\lambda_2=\lambda_2^{k-1}$、$\lambda_3=\lambda_3^{k-1}$。

\textbf{步骤2：}重复以下过程。
\begin{enumerate}
\item 令
\[
Z=\begin{bmatrix}Z_\Gamma\\Z_A\\Z_B\end{bmatrix}
=\begin{bmatrix}
\operatorname{Prox}_{\lambda_1\Omega_\Gamma}(u_\Gamma^k)\\
\operatorname{Prox}_{\lambda_2\Omega_A}(u_A^k)\\
\operatorname{Prox}_{\lambda_3\Omega_B}(u_B^k)
\end{bmatrix}.
\]
逐元素计算近端算子：
\[
\begin{aligned}
u_\Gamma^k&=W_\Gamma^k-\lambda_1\left.\frac{\partial f}{\partial W_\Gamma}\right|_{W^k},\\
u_A^k&=W_A^k-\lambda_2\left.\frac{\partial f}{\partial W_A}\right|_{W^k},\qquad
u_B^k=W_B^k-\lambda_3\left.\frac{\partial f}{\partial W_B}\right|_{W^k},
\end{aligned}
\]
\[
\begin{gathered}
W_\Gamma^k=(\gamma_{ij}^k)_{MT\times MT},\quad
W_A^k=(\alpha_{ij}^k)_{MT\times MT},\quad W_B^k=(\beta_{ij}^k)_{KT\times MT},\\
u_\Gamma^k=(u_{\Gamma ij}^k)_{MT\times MT},\quad
u_A^k=(u_{Aij}^k)_{MT\times MT},\quad u_B^k=(u_{Bij})_{KT\times MT},
\end{gathered}
\]
\begin{equation}
\gamma_{ij}^z=\operatorname{sign}(u_{\Gamma ij}^k)(|u_{\Gamma ij}^k|-\lambda_1^k)_+,
\quad i,j=1,\ldots,MT,\tag{4.G.23}
\end{equation}
\begin{equation}
\alpha_{ij}^z=\operatorname{sign}(u_{Aij}^k)(|u_{Aij}^k|-\lambda_2^k)_+,
\quad i,j=1,\ldots,MT,\tag{4.G.24}
\end{equation}
\begin{equation}
\beta_{ij}^z=\operatorname{sign}(u_{Bij}^k)(|u_{Bij}^k|-\lambda_3^k)_+,
\quad i=1,\ldots,KT,\ j=1,\ldots,MT,\tag{4.G.25}
\end{equation}
\begin{equation}
\Gamma_z=(\gamma_{ij}^z)_{MT\times MT},\quad
A_z=(\alpha_{ij}^z)_{MT\times MT},\quad B_z=(\beta_{ij})_{KT\times MT}.\tag{4.G.26}
\end{equation}
\[
Z_\Gamma=\operatorname{vec}(\Gamma_z),\quad
Z_A=\operatorname{vec}(A_z),\quad Z_B=\operatorname{vec}(B_z).
\]
\item 计算辅助函数：
\[
\begin{aligned}
f_\lambda(Z,G^k)={}&f(Z)
+\left(\frac{\partial f}{\partial Z}\right)^T(G^k-Z)
+\frac1{2\lambda_1}\|G_\Gamma^k-Z_\Gamma\|_2^2\\
&{}+\frac1{2\lambda_2}\|G_A^k-Z_A\|_2^2
+\frac1{2\lambda_3}\|G_B^k-Z_B\|_2^2.
\end{aligned}
\]
\item 若 $f(Z)\le\widehat f_\lambda(Z,G^k)$，则退出循环，转至步骤3。
\item 更新 $\lambda_1=\alpha_1\lambda_1$、$\lambda_2=\alpha_2\lambda_2$、$\lambda_3=\alpha_3\lambda_3$。
\end{enumerate}
\textbf{步骤3：}令 $\lambda_1^k=\lambda_1$、$\lambda_2^k=\lambda_2$、$\lambda_3^k=\lambda_3$，
$G^{k+1}=Z$，以及
\begin{equation}
W_\Gamma^{k+1}=\Gamma_Z,\quad W_A^{k+1}=A_Z,\quad
W_B^{k+1}=B_Z,\quad k=k+1.\tag{4.G.27}
\end{equation}
\end{bookalgorithm}

最后，矩阵初值可由最小二乘估计得到。令
\[
\begin{gathered}
W^t=\begin{bmatrix}Y_l^{t-1}&X\end{bmatrix},\quad
W=\begin{bmatrix}W^1&\cdots&W^T\end{bmatrix},\quad
Z_i^t=\begin{bmatrix}Y_i^t&Y_{li}^{t-1}&X_i\end{bmatrix},\quad
\Lambda_i=W'W;
\end{gathered}
\]
\[
y_i=\begin{bmatrix}y_i^1\\y_i^2\\\vdots\\y_i^T\end{bmatrix},\quad
Z_i=\begin{bmatrix}Z_i^1&0&\cdots&0\\0&Z_i^2&\cdots&0\\\vdots&\vdots&\vdots&\vdots\\0&0&\cdots&Z_i^T\end{bmatrix},\quad
\delta_i=\begin{bmatrix}\delta_i^1\\\delta_i^2\\\vdots\\\delta_i^T\end{bmatrix},\quad
\delta_i^t=\begin{bmatrix}\gamma_i^t\\\alpha_i^t\\\beta_i^t\end{bmatrix}.
\]
参数 $\delta_i$ 的最小二乘估计量为
\begin{equation}
\widehat\delta_i=[(W'Z_i)'\Lambda_i^{-1}W'Z_i]^{-1}
(W'Z_i)'\Lambda_i^{-1}W'y_i,\quad i=1,\ldots,M.\tag{4.G.28}
\end{equation}
从 $\widehat\delta_i$ 提取 $\gamma_i^t,\alpha_i^t,\beta_i^t$，并组装为
\[
\Gamma^t=\begin{bmatrix}\gamma_1^t&\cdots&\gamma_M^t\end{bmatrix},\quad
A^t=\begin{bmatrix}\alpha_1^t&\cdots&\alpha_M^t\end{bmatrix},\quad
B^t=\begin{bmatrix}\beta_1^t&\cdots&\beta_M^t\end{bmatrix}.
\]
定义
\[
\Gamma_0=\begin{bmatrix}\Gamma^1&0&\cdots&0\\0&\Gamma^2&\cdots&0\\\vdots&\vdots&\vdots&\vdots\\0&0&\cdots&\Gamma^T\end{bmatrix},\quad
A_0=\begin{bmatrix}A^1&0&\cdots&0\\0&A^2&\cdots&0\\\vdots&\vdots&\vdots&\vdots\\0&0&\cdots&A^T\end{bmatrix},
\]
\[
B_0=\begin{bmatrix}B^1&0&\cdots&0\\0&B^2&\cdots&0\\\vdots&\vdots&\vdots&\vdots\\0&0&\cdots&B^T\end{bmatrix}.
\]
设置矩阵 $W^0$ 与向量 $G^0$ 的初值：
\[
W_\Gamma^0=\Gamma_0,\qquad W_A^{01}=A_0,\qquad W_B^0=B_0,
\]
\[
G_\Gamma^0=\operatorname{vec}(\Gamma_0),\quad
G_A^0=\operatorname{vec}(A_0),\quad G_B^0=\operatorname{vec}(B_0),\qquad
G^0\begin{bmatrix}G_\Gamma^0\\G_A^0\\G_B^0\end{bmatrix}.
\]
